\documentclass[14pt,dvipdfmx]{jsarticle}
\input{lecture-preamble.tex}

\setcounter{definition}{2}
\setcounter{theorem}{7}
\setcounter{technique}{2}

\begin{document}
\lectureheader{数学Ⅱ}{5}{指数関数と対数関数}{2}{対数関数}{5-A}{常用対数}{169}{常用対数の意味を理解し, 常用対数表を使って値を求めよう}

% ============ 1ページ目（表・左） ============
\vspace{-1.6zw}
\heading{〜前時の復習〜}
\vspace{-0.35em}
\begin{reviewbox}{対数関数を含む方程式・不等式}
  \begin{enumerate}[label=\textbf{\arabic*}, leftmargin=1.5zw, itemsep=0.02em]
    \item 真数が\fitblankbf{\text{正}}である範囲（\fitblankbf{\text{真数条件}}）を求める.
    \item 対数の定義や, 対数の性質を使って \fitblankbf{式を変形}する.
    \item 対数不等式の底が $a>1$ なら不等号の向きは\fitblankbf{\text{そのまま}},
    $0<a<1$ なら不等号の向きは\fitblankbf{\text{逆}}にする.
    \item 真数条件との \fitblankbf{\text{吟味}}をする.
  \end{enumerate}
\end{reviewbox}

\vspace{0.25em}
地球の質量、太陽の質量はそれぞれ
\par\vspace{0.2em}
{\centering
約 $5,970,000,000,000,000,000,000,000$ kg,

約 $1,990,000,000,000,000,000,000,000,000,000$ kg
\par}
であることが知られている.これは直感的にわかりにくいため、指数を使って表記してみると、それぞれ

\par\vspace{0.2em}
{\centering
約 $\fitblankbf{5.97 \times 10^{24}}$ kg, \quad
約 $\fitblankbf{1.99 \times 10^{30}}$ kg
\par}

と表せる.ここで、$10$ を底とする対数を考え、対数の性質を使って変形すると、
\vspace{0.2em}
\begin{answerbox}{}[14]
\begin{align*}
  \vspace{-0.5em}
  \log_{10} (5.97 \times 10^{24})
  &= \log_{10} 5.97 + \log_{10} 10^{24} \\
  &= \log_{10} 5.97 + 24 \\
  &= 0.7760 + 24 \\
  &= 24.7760
\end{align*}
$\therefore \text{(地球の質量)} = 10^{24.7760}(\text{kg})$
\end{answerbox}

\vfill
\nextpage

% ============ 2ページ目（表・右） ============

\begin{answerbox}{}[14]
  \begin{align*}
    \vspace{-1em}
    \log_{10} (1.99 \times 10^{30}) &= \log_{10} 1.99 + \log_{10} 10^{30} \\
    &= \log_{10} 1.99 + 30 \\
    &= 0.2989 + 30 \\
    &= 30.2989
  \end{align*}
  $\therefore \text{(太陽の質量)} = 10^{30.2989}(\text{kg})$
\end{answerbox}
\vspace{1em}

この計算にはいろいろなメリットがある。非常に大きな累乗の計算や、非常に小さな累乗の計算が容易にできる。

また、非常に大きな数や非常に小さな数の \fitblankbf{桁数} を求めることができる。

例えば、地球の質量は $10^{24.7760}$ kg であるから、\fitblankbf{25桁}の数であると分かる。
\vspace{2em}

\heading{常用対数}

\begin{definitionbox}{常用対数}
  $10$ を底とする対数を, \fitblankbf{\text{常用対数}}という.
\end{definitionbox}

\vspace{0.15em}

\vspace{0.08em}
\textbf{注意}\quad $1 \leqq a < 10$ であるから, $\log_{10} a$ の値の範囲は
$0 \leqq \log_{10} a < 1$ である.

\vspace{0.25em}

\vfill
\nextpage

% ============ 3ページ目（裏・左） ============

\begin{exampleprob}{12}
  常用対数表を用いて, 次の値を求めよ。
  \begin{tasks}[label=(\arabic*), label-width=1.6zw, item-indent=2zw, after-item-skip=0.08em](2)
    \task $\log_{10} 1620000$
    \task $\log_{10} 0.00162$
  \end{tasks}
\end{exampleprob}

\vspace{-1.4zw}
\begin{answerbox}{(1)}[12]
  $1620000 = 1.62 \times 10^{6}$ であるから\\
  $\begin{aligned}[t]
    \log_{10} 1620000 &= \log_{10} 1.62 + \log_{10} 10^{6} \\
    &= 0.2095 + 6 \\
    &= 6.2095
  \end{aligned}$
\end{answerbox}
\vspace{0.12em}
\begin{answerbox}{(2)}[12]
  $0.00162 = 1.62 \times 10^{-3}$ であるから\\
  $\begin{aligned}[t]
    \log_{10} 0.00162 &= \log_{10} 1.62 + \log_{10} 10^{-3} \\
    &= 0.2095 - 3 \\
    &= -2.7905
  \end{aligned}$
\end{answerbox}

\vfill
\nextpage

% ============ 4ページ目（裏・右） ============
\begin{practiceprob}{24}
  常用対数表を用いて, 次の値を小数第$4$位まで求めよ。
  \begin{tasks}[label=(\arabic*), label-width=1.6zw, item-indent=2zw, after-item-skip=0.06em](3)
    \task $\log_{10} 3450$
    \task $\log_{10} 92000$
    \task $\log_{10} 0.000618$
  \end{tasks}
\end{practiceprob}

\begin{answerbox}{(1)}[12]
  $3450 = 3.45 \times 10^{3}$ であるから\\
  $\begin{aligned}[t]
    \log_{10} 3450 &= \log_{10} 3.45 + \log_{10} 10^{3} \\
    &= 0.5378 + 3 \\
    &= 3.5378
  \end{aligned}$
\end{answerbox}
\vspace{0.08em}
\begin{answerbox}{(2)}[12]
  $92000 = 9.20 \times 10^{4}$ であるから\\
  $\begin{aligned}[t]
    \log_{10} 92000 &= \log_{10} 9.20 + \log_{10} 10^{4} \\
    &= 0.9638 + 4 \\
    &= 4.9638
  \end{aligned}$
\end{answerbox}
\vspace{0.08em}
\begin{answerbox}{(3)}[12]
  $0.000618 = 6.18 \times 10^{-4}$ であるから\\
  $\begin{aligned}[t]
    \log_{10} 0.000618 &= \log_{10} 6.18 + \log_{10} 10^{-4} \\
    &= 0.7910 - 4 \\
    &= -3.2090
  \end{aligned}$
\end{answerbox}

\vspace{0.35em}
\begin{summarybox}
  \begin{enumerate}[label=\textbf{\arabic*}, leftmargin=1.5zw, itemsep=0.08em]
    \item $10$ を底とする対数を\fitblankbf{\text{常用対数}}という.
  \end{enumerate}
\end{summarybox}

\vfill
\nexttime{常用対数を使って, 桁数を求めてみよう!}

\end{document}
